% Source: https://artofproblemsolving.com/community/c3h1646397_an_algebra_question
% Stable identity: aops-topic-1646397-q01
\textbf{(a)} Let the roots be $r_1$ and $r_2$ (with multiplicity). By Vieta's formulas, $r_1+r_2=-\dfrac {b}{a}$, which is real when $a,b$ are real. If $r_1$ is real, then $r_2=-\dfrac {b}{a}-r_1$ is real too. So with real coefficients either both roots are real or neither is. (Equivalently: non-real roots of real polynomials come in conjugate pairs.)

\textbf{(b)} The equation $x^2-ix=0$, that is $x(x-i)=0$, has the real root $0$ and the non-real root $i$.

So the answer is $\boxed{\text{no for real coefficients; yes in general, e.g. } x^2-ix=0}$.
